\documentclass[dvipdfmx, uplatex]{jsarticle}
\usepackage[dvipdfmx]{graphicx}
\usepackage{amsmath}
\usepackage{amssymb}
\usepackage{amsfonts}
\usepackage{amsthm}
\usepackage{mathrsfs}
\usepackage{bm}
\usepackage{braket}
\usepackage{xr}
\externaldocument{exercise_so2_index_theorem}

\usepackage{silence}
\WarningFilter*{hyperref}{Token not allowed in a PDF string}
\usepackage[dvipdfmx]{hyperref}
\usepackage{pxjahyper}
\allowdisplaybreaks[4]

\usepackage[left=25truemm,top=20truemm,bottom=20truemm,right=25truemm]{geometry}

\newcommand{\zbar}{\bar z}
\newcommand{\dz}{\partial_z}
\newcommand{\dzb}{\partial_{\bar z}}

\theoremstyle{definition}
\newtheorem*{sol}{解答}

\title{解答：回転群 $SO(2)$ の表現から指数定理へ\\
\large ―― \texttt{exercise\_so2\_index\_theorem.tex} の解答 ――}
\author{}
\date{}

\begin{document}
\maketitle
\vspace{-6zh}

記号・問題番号は問題ファイル \texttt{exercise\_so2\_index\_theorem.tex} のものを踏襲する。
\textbf{各問の解（採点用）}は太字（\boxed）で示す。

%======================================================================
\section{SO(2) の二つの表現}
%======================================================================
\begin{sol}
\textbf{(1) ベクトル表現の群則。}加法定理より
\[
    R(\theta_1)R(\theta_2)=
    \begin{pmatrix}c_1c_2-s_1s_2 & -(s_1c_2+c_1s_2)\\ s_1c_2+c_1s_2 & c_1c_2-s_1s_2\end{pmatrix}
    =\begin{pmatrix}\cos(\theta_1{+}\theta_2)&-\sin(\theta_1{+}\theta_2)\\ \sin(\theta_1{+}\theta_2)&\cos(\theta_1{+}\theta_2)\end{pmatrix}
    =\boxed{R(\theta_1+\theta_2)}
\]
（$c_i=\cos\theta_i,\ s_i=\sin\theta_i$）。$\boxed{R(0)=I}$，$\boxed{R(\theta)^{-1}=R(-\theta)=R(\theta)^{\mathsf T}}$。

\par\smallskip
\textbf{(2) 生成子の作用。}$\boxed{Gx=y,\quad Gy=-x}$（$Gx=y\partial_x x-x\partial_y x=y$ 等）。
$u=e^{\theta G}x,\ v=e^{\theta G}y$ は $u'=e^{\theta G}Gx=v,\ v'=e^{\theta G}Gy=-u$ を満たす。
$u''=-u$，初期条件 $u(0)=x,v(0)=y$ より
\[
    \boxed{\,e^{\theta G}x=x\cos\theta+y\sin\theta,\qquad e^{\theta G}y=-x\sin\theta+y\cos\theta\,}.
\]

\par\smallskip
\textbf{(3) 関数表現の群則。}$[G,G]=0$ ゆえ $e^{\theta_1G}e^{\theta_2G}=e^{(\theta_1+\theta_2)G}$，すなわち
$\boxed{U(\theta_1)U(\theta_2)=U(\theta_1+\theta_2)}$。$\boxed{U(0)=\mathrm{id},\ U(\theta)^{-1}=U(-\theta)}$。
1 次関数 $x,y$ への作用は (2) より
$U(\theta)\begin{pmatrix}x\\ y\end{pmatrix}=\begin{pmatrix}\cos\theta & \sin\theta\\ -\sin\theta & \cos\theta\end{pmatrix}
\begin{pmatrix}x\\ y\end{pmatrix}$ ゆえ，$\mathrm{span}\{x,y\}$ 上で
$\boxed{U(\theta)\big|_{\{x,y\}}=R(-\theta)=R(\theta)^{\mathsf T}}$（反傾＝双対表現；$SO(2)$ では
ベクトル表現と同値）。よって\emph{同じ回転群}が，2 次元ベクトル表現としても無限次元関数表現としても
実現され，後者の 1 次部分に前者が埋め込まれている。
\end{sol}

%======================================================================
\section{固有ベクトルと固有関数}
%======================================================================
\begin{sol}
\textbf{(4) ベクトル表現の固有ベクトル。}特性方程式 $(\cos\theta-\lambda)^2+\sin^2\theta=0$ より
$\boxed{\lambda_\pm(\theta)=e^{\pm i\theta}}$。固有ベクトルは
\[
    \boxed{\,\bm e_+=\begin{pmatrix}1\\ -i\end{pmatrix}\ (\lambda_+=e^{i\theta}),\qquad
    \bm e_-=\begin{pmatrix}1\\ i\end{pmatrix}\ (\lambda_-=e^{-i\theta})\,}
\]
（$R(\theta)\bm e_\pm=e^{\pm i\theta}\bm e_\pm$，$\theta$ によらない定ベクトル）。\textbf{2 個}。

\par\smallskip
\textbf{(5) 関数表現の固有関数。}$\partial_\varphi=-y\partial_x+x\partial_y=-G$ ゆえ $\boxed{G=-\partial_\varphi}$。
$G\psi=\mu\psi\Leftrightarrow-\partial_\varphi\psi=\mu\psi\Leftrightarrow\psi\propto e^{-\mu\varphi}$。
単価性より $-\mu=in$（$n\in\mathbb{Z}$），すなわち
\[
    \boxed{\,\psi_n=e^{in\varphi}\times(\text{任意の動径関数}),\quad G\psi_n=-in\,\psi_n,\quad
    U(\theta)\psi_n=e^{-in\theta}\psi_n\,}.
\]
固有関数は $\mathbb{Z}$ で添字づけられ\textbf{可算無限個}。ベクトル表現の固有値 $e^{\pm i\theta}$ は
$n=\mp1$ の場合に当たる（ベクトル表現 $=$ 関数表現の $n=\pm1$ 部分表現）。
\end{sol}

%======================================================================
\section{複素座標と Laurent 基底}
%======================================================================
\begin{sol}
\textbf{(6) 複素座標での生成子。}$Gz=Gx+iGy=y-ix=-i(x+iy)$，$G\bar z=y+ix=i(x-iy)$ ゆえ
$\boxed{Gz=-iz,\quad G\bar z=+i\bar z}$。これより
\[
    \boxed{\,G=-i\bigl(z\partial_z-\bar z\partial_{\bar z}\bigr)\,}
\]
（$z\partial_z,\bar z\partial_{\bar z}$ への作用が一致することで確認できる）。単項式について
$G(z^n\bar z^m)=-i(n-m)\,z^n\bar z^m$ ゆえ固有関数であり，
\[
    \boxed{\,G\text{ 固有値}=-i(n-m),\qquad U(\theta)\text{ 固有値}=e^{-i(n-m)\theta}\,}.
\]

\par\smallskip
\textbf{(7) Laurent 基底。}単位円 $|z|=1$ では $z=e^{i\varphi}$ ゆえ (5) の固有関数は
$\boxed{e^{in\varphi}=z^n}$。$\{z^n\}_{n\in\mathbb{Z}}$ は環状領域（あるいは $|z|=1$）上の関数の
\textbf{Laurent 級数の基底}である（$n\ge0$ が正則部分 $z^n$，$n<0$ が主要部 $z^{-|n|}$）。
回転生成子の固有関数がそのまま Laurent 基底になっている。
\end{sol}

%======================================================================
\section{Poincar\'e--Hopf：指数の総和}
%======================================================================
\begin{sol}
\textbf{(8) 指数の総和。}有限平面では零点の位数 $=+$ 指数，極の位数 $=-$ 指数。無限遠は
$v=-w^2R(1/w)\partial_w$ の $w=0$ での位数で読む。
\[
\renewcommand{\arraystretch}{1.3}
\begin{array}{c|c|c|c}
 & \text{有限の零点／極（指数）} & \text{無限遠 }w=0\ (\text{指数}) & \text{総和}\\\hline
\text{(a) }\partial_z & \text{なし} & -w^2\cdot1=-w^2\partial_w\ (+2) & \boxed{+2}\\
\text{(b) }z\partial_z & z=0\ (+1) & -w^2\cdot\tfrac1w=-w\partial_w\ (+1) & \boxed{+2}\\
\text{(c) }\tfrac1z\partial_z & z=0\ \text{極}\ (-1) & -w^2\cdot w=-w^3\partial_w\ (+3) & \boxed{+2}\\
\text{(d) }(z^2-1)\partial_z & z=\pm1\ (+1,+1) & -(1-w^2)\partial_w\ (\text{零点なし}) & \boxed{+2}
\end{array}
\]
いずれも総和 $\boxed{2}$。これは $\boxed{S^2\text{ の Euler 数 }\chi(S^2)=2}$ に等しい（\textbf{Poincar\'e--Hopf
定理}）。一般に種数 $g$ の閉曲面上の有理型ベクトル場では（零点数）$-$（極数）$=\chi=2-2g$。
(b) の零点 $z=0,\infty$ は (2)(4) の回転の固定点（北極・南極）に対応し，回転ベクトル場が指数
$+1$ の固定点を 2 個持つことが $\chi=2$ を与える。
\end{sol}

%======================================================================
\section{U(1) 接続・Cauchy--Riemann 作用素・Chern 数}
%======================================================================
\begin{sol}
\textbf{(9) ゲージ共変性。}$D_{\bar z}(e^{i\alpha}\psi)=e^{i\alpha}\bigl(i(\dzb\alpha)\psi+\dzb\psi+iA'_{\bar z}\psi\bigr)$
が $e^{i\alpha}(\dzb\psi+iA_{\bar z}\psi)$ に等しいためには
\[
    \boxed{\,A'_{\bar z}=A_{\bar z}-\dzb\alpha\,}\qquad(\text{1 形式では }A\to A-d\alpha).
\]
曲率 $F=dA\to d(A-d\alpha)=dA=F$ ゆえ\textbf{ゲージ不変}（$d^2=0$）。

\par\smallskip
\textbf{(10) 単極子と Chern 数。}
\[
    \boxed{\,F=dA=\frac{n}{2}\sin\vartheta\,d\vartheta\wedge d\varphi\,},\qquad
    \int_{S^2}F=\frac n2\int_0^\pi\!\sin\vartheta\,d\vartheta\int_0^{2\pi}\!d\varphi=\frac n2\cdot2\cdot2\pi=\boxed{2\pi n}.
\]
よって第一 Chern 数 $\boxed{c_1=\dfrac{1}{2\pi}\displaystyle\int_{S^2}F=n\in\mathbb{Z}}$（Dirac の量子化）。
（立体射影では $F=n\,\dfrac{i\,dz\wedge d\bar z}{(1+|z|^2)^2}$ と同じものになり，
$\int i\,dz\wedge d\bar z/(1+|z|^2)^2=2\pi$ から同じ結論を得る。）

\par\smallskip
\textbf{(11) 正則ゼロモードと閉曲面の指数。}正則ゲージで $D_{\bar z}\psi=0\Leftrightarrow\psi=f(z)\times$（正則枠）。
$S^2$ 上 2 乗可積分（南極 $w=1/z$ でも正則）となるのは $f$ が高々 $n$ 次の多項式，すなわち
$f=1,z,\dots,z^n$。よって
\[
    \boxed{\dim\ker D_{\bar z}=n+1},\qquad
    \boxed{\mathrm{ind}\,D_{\bar z}=(n+1)-0=n+1}.
\]
これは $\mathrm{ind}=c_1+\tfrac12\chi(S^2)=n+\tfrac12\cdot2=n+1$（Atiyah--Singer／Riemann--Roch，
$\mathrm{ind}=\deg+1-g$，$g=0$）と一致する。ゼロモード $z^k$ は第 6 節の Laurent 基底の
$0\le k\le n$ 部分にほかならない。
\end{sol}

%======================================================================
\section{半球分割・境界作用素・APS 指数定理}
%======================================================================
\begin{sol}
\textbf{(12) 赤道上の境界作用素。}最低 Landau 準位は全角運動量 $J=n/2$ の多重項で，
$\psi_k\propto z^k(1+|z|^2)^{-n/2}$ の $J_z$ 固有値は
\[
    \boxed{\,J_z\,\psi_k=\Bigl(k-\tfrac n2\Bigr)\psi_k,\qquad k=0,1,\dots,n\,}
    \quad(J_z=-\tfrac n2,\dots,+\tfrac n2).
\]
赤道 $|z|=1$ 上で $\psi_k\propto e^{ik\varphi}$ ゆえ，境界に誘導される 1 次元作用素 $B$（軌道部分は
回転生成子 $iG=-i\partial_\varphi$ に単極子の寄与を加えた全角運動量 $J_z$）の固有関数は\textbf{Laurent 基底
$z^k=e^{ik\varphi}$}，固有値は $\boxed{k-\tfrac n2}$ である。回転の生成子が境界作用素として自然に現れている。

\par\smallskip
\textbf{(13) $\eta$ 不変量。}スペクトル $\{m+c\}$（固有関数 $e^{im\varphi}$，$0<c<1$）について
$\eta(s)=\zeta(s,c)-\zeta(s,1-c)$，$\zeta(0,a)=\tfrac12-a$ より
\[
    \boxed{\,\eta(0)=\bigl(\tfrac12-c\bigr)-\bigl(\tfrac12-(1-c)\bigr)=1-2c\,}\qquad(0<c<1).
\]
単極子では $c\equiv\tfrac n2\pmod1$：
\begin{itemize}
\item[・]\textbf{$n$ 偶数}（$c\equiv0$）：スペクトル $\{m\}$ は $0$ に関して対称ゆえ $\boxed{\eta(0)=0}$，
かつ零固有値（$k=n/2$）が 1 個 $\boxed{h=1}$。
\item[・]\textbf{$n$ 奇数}（$c\equiv\tfrac12$）：スペクトル $\{m+\tfrac12\}$ も対称ゆえ $\boxed{\eta(0)=0}$，
零固有値なし $\boxed{h=0}$。
\end{itemize}
（いずれもスペクトルが $0$ に関して対称なので $\eta(0)=0$。）

\par\smallskip
\textbf{(14) APS の検証。}北半球の APS 境界条件は\emph{境界固有値 $k-\tfrac n2<0$ のモードを残す}。
生き残る正則切断は $z^k$（$0\le k<\tfrac n2$）で $N_+=\#\{k:0\le k<\tfrac n2\}=\lceil n/2\rceil$，
南半球は $k>\tfrac n2$ を残し $N_-=n+1-\lceil n/2\rceil$（$n$ 偶数の境界零モード $k=\tfrac n2$ は $h/2$ で配分）。
\[
\renewcommand{\arraystretch}{1.35}
\begin{array}{c|c|c|c|c|c|c}
n & c_1 & \mathrm{ind}_{S^2} & \{J_z\}=\{k-\tfrac n2\} & N_+ & N_- & N_++N_-\\\hline
0 & 0 & 1 & \{0\} & 0 & 0\,(+\text{境界零モード}) & \boxed{1}\\
1 & 1 & 2 & \{-\tfrac12,+\tfrac12\} & 1 & 1 & \boxed{2}\\
2 & 2 & 3 & \{-1,0,+1\} & 1 & 1\,(+\text{境界零モード}) & \boxed{3}
\end{array}
\]
各行で $\boxed{N_++N_-=\mathrm{ind}_{S^2}=n+1}$（境界 $J_z=0$ のモードは両半球に $h/2$ ずつ配分）。
n+1 個の正則切断が，赤道上の $J_z$ の符号で南北に振り分けられている。

\smallskip
APS 指数定理の右辺：北半球のバルクは
$\tfrac{1}{2\pi}\int_{D_+}F=\tfrac n2$（対称ゆえ半分）と
$\tfrac12\cdot\tfrac{1}{2\pi}\int_{D_+}R_{\rm Euler}=\tfrac12\cdot1=\tfrac12$（赤道は測地線で測地的曲率 $0$，
Gauss--Bonnet より $\tfrac{1}{2\pi}\int_{D_+}R_{\rm Euler}=\chi(D_+)=1$）の和 $\tfrac n2+\tfrac12$。よって
\[
    \mathrm{ind}(D_{\bar z},\text{APS})=\Bigl(\tfrac n2+\tfrac12\Bigr)-\frac{\eta(0)+h}{2}
    =\begin{cases}\tfrac n2+\tfrac12-\tfrac{0+1}{2}=\dfrac n2 & (n\ \text{偶}),\\[1.2ex]
    \tfrac n2+\tfrac12-\tfrac{0+0}{2}=\dfrac{n+1}{2} & (n\ \text{奇}),\end{cases}
    =\boxed{\Bigl\lceil\tfrac n2\Bigr\rceil=N_+}.
\]
具体的には $n=0\!:0,\ n=1\!:1,\ n=2\!:1$ で前表の $N_+$ に一致する。バルクの Chern 数（${+}$Euler 項）と
境界のスペクトル非対称性 $\eta$ の差として半球の指数が定まり，南北を貼り合わせると境界項が相殺して
閉曲面の指数 $n+1=c_1+\tfrac12\chi$ を再現する――これが APS 指数定理の主張である。

\par\smallskip
\noindent\textbf{参考。}指数定理：M.\,F.~Atiyah, I.\,M.~Singer, \textit{The index of elliptic operators}
(Ann.\ Math.\ 1968)。境界付き版：M.\,F.~Atiyah, V.\,K.~Patodi, I.\,M.~Singer,
\textit{Spectral asymmetry and Riemannian geometry I} (Math.\ Proc.\ Camb.\ Phil.\ Soc.\ 77 (1975) 43)。
球面上の Landau 準位・単極子調和関数：F.\,D.\,M.~Haldane (Phys.\ Rev.\ Lett.\ 51 (1983) 605)，
T.\,T.~Wu, C.\,N.~Yang (Nucl.\ Phys.\ B107 (1976) 365)。
\end{sol}

\end{document}
